% Desarrollo separado, 9 de septiembre de 2026.
% Amplía el lector producto de centro_electronico_registros.tex.
% No identifica la media vuelta espacial con una conjugación física.
\subsubsection{Bidirectionnalité du lecteur produit et mémoire des quotients}
\label{bim09:seccion}

La lecture du produit conserve quatre coordonnées arithmétiques avant de déterminer une signature. Pour $x=(i,j)\in D_9^2$, $D_9=\{1,\ldots,9\}$, on définit
\[
 N(x)=ij,\qquad r(x)=1+((ij-1)\bmod9),\qquad
 q(x)=\frac{N(x)-r(x)}9,
\]
et le lecteur discret
\[
 \ell(1,2,4,8,7,5)=(0,1,2,3,4,5),\qquad
 \ell(3)=\ell(6)=\ell(9)=0.
\]
Ainsi, $N=r+9q$ conserve simultanément le produit entier, le résidu positif
et le quotient arithmétique. La fonction $f_\ell(x)=\ell(r(x))$ est la
coordonnée utilisée par la signature régionale. Les fonctions $N,r,q$ et
$f_\ell$ ont un domaine positionnel ; la direction de transport demeure
dans le germe orienté.

\paragraph{Avancée, demi-tour et lecture avant le transport.}

Soit $\mathscr X_\times=D_9^2\times\{N,E,S,O\}$. Les vecteurs cardinaux
sont $\nu_N=(-1,0)$, $\nu_E=(0,1)$, $\nu_S=(1,0)$,
$\nu_O=(0,-1)$. La somme positionnelle est réduite aux représentants positifs
modulo neuf. Définissons
\[
 P(x,d)=(x+\nu_d,d),\qquad H(x,d)=(x,\bar d),
 \qquad \nu_{\bar d}=-\nu_d.
\]
Ces opérateurs satisfont $P^9=I$, $H^2=I$ et $HPH=P^{-1}$.

Dans les cinquante-quatre premières transitions TPK, le canal produit
est activé aux phases $2,5,8$ de chaque fenêtre nonadique. Chaque activation
lit la position avant l'avancée. Après les transitions $27$ et $54$,
on applique $H$. Les dix-huit positions lues sont donc, pour un
germe $s$,
\[
 P^0s,P^1s,\ldots,P^8s,\quad
 HP^9s,\,P H P^9s,\ldots,P^8H P^9s,
\]
étant entendu que la lecture positionnelle omet uniquement la direction.
Pour toute fonction positionnelle $f$, posons $f_n=f(P^ns)$, avec
indices modulo neuf. Le mot des lectures brutes est
\begin{equation}
 \mathcal R_f(s)=
 (f_0,f_1,\ldots,f_8,\ f_0,f_8,f_7,\ldots,f_1).
 \label{bim09:palabra}
\end{equation}
Les six sommes de fenêtre, comprenant chacune trois activations, sont
\begin{align}
 A_j^f(s)&=f_{3j}+f_{3j+1}+f_{3j+2},&&0\leq j<3,\label{bim09:ida}\\
 A_{3+j}^f(s)&=f_{-3j}+f_{-3j-1}+f_{-3j-2},&&0\leq j<3.
 \label{bim09:vuelta}
\end{align}
Pour $f=f_\ell$, la signature est
$E_{\rm sig}(s)=([A_0^f]_{10},\ldots,[A_5^f]_{10})$.
Le quotient décimal de chaque accumulateur est
$c_j^{10}=\lfloor A_j^f/10\rfloor$ et conserve
$A_j^f=[A_j^f]_{10}+10c_j^{10}$.

\begin{proposition}[Demi-tour et permutation des demi-cycles]
\label{bim09:media-vuelta}
Pour tout germe orienté et toute fonction positionnelle $f$,
\[
 \mathcal R_f(Hs)=\operatorname{rot}_9\mathcal R_f(s),\qquad
 A_j^f(Hs)=A_{j+3\bmod6}^f(s).
\]
La même permutation transporte les résidus et les quotients décimaux
des accumulateurs. Elle s'applique en particulier à $N,r,q$ et $f_\ell$.
\end{proposition}
\begin{proof}
La relation $P^nH=HP^{-n}$ et l'égalité $f(Hs)=f(s)$ donnent
$f(P^nHs)=f(P^{-n}s)=f_{-n}$. Substituer dans
\eqref{bim09:palabra} échange ses deux groupes de neuf entrées.
La sommation par groupes de trois produit l'identité des fenêtres. La division
euclidienne par dix s'effectue composante par composante et respecte la
permutation.
\end{proof}

L'avancée possède également une loi exacte sur les accumulateurs :
\begin{align}
 A_j^f(Ps)-A_j^f(s)&=f_{3j+3}-f_{3j},&&0\leq j<3,
 \label{bim09:avance-ida}\\
 A_{3+j}^f(Ps)-A_{3+j}^f(s)&=f_{1-3j}-f_{-2-3j},&&0\leq j<3.
 \label{bim09:avance-vuelta}
\end{align}
La démonstration consiste à annuler les deux termes communs de chaque fenêtre. Le transport d'une signature exige donc les évaluations de frontière qui apparaissent dans ces différences, en plus des sommes déjà exprimées.

\paragraph{Inversion chronologique et termes de frontière.}

Soit $\gamma=(x_0,\ldots,x_n)$ une trajectoire enregistrée et soit
$\gamma^\dagger=(x_n,\ldots,x_0)$ son inversion chronologique, chaque
arête étant orientée en sens opposé. Pour une fonction positionnelle $f$,
la lecture précédant chaque avancée est
\[
 C_f(\gamma)=\sum_{t=0}^{n-1}f(x_t).
\]
Alors
\begin{equation}
 C_f(\gamma^\dagger)=C_f(\gamma)+f(x_n)-f(x_0).
 \label{bim09:frontera-inversion}
\end{equation}
Si $n=3m$ et $A_j^f$ regroupe trois arêtes consécutives,
\begin{equation}
 A_j^f(\gamma^\dagger)=A_{m-1-j}^f(\gamma)
 +f(x_{3(m-j)})-f(x_{3(m-j-1)}).
 \label{bim09:frontera-ventana}
\end{equation}
Les deux formules résultent du fait que la trajectoire inverse lit les extrémités
d'arrivée des arêtes originales. Les termes intérieurs coïncident ;
la différence conserve les extrémités.

Pour un observable d'arête antisymétrique $a(x,y)=-a(y,x)$, son mot
se transforme effectivement selon
\[
 \mathcal R_a(\gamma^\dagger)
 =-\operatorname{rev}\mathcal R_a(\gamma).
\]
Cette identité agit sur les arêtes orientées. Les identités
\eqref{bim09:frontera-inversion}--\eqref{bim09:frontera-ventana}
agissent sur les lectures positionnelles. Le demi-tour $H$ change la
direction initiale et conserve l'ordre du calendrier ; l'inversion
chronologique inverse le registre complet. Leurs domaines et leurs règles de
lecture sont spécifiés séparément.

\paragraph{Cocycle d'occupation et mémoire conservée à l'aller et au retour.}

Pour une arête cardinale $x\to y$, de représentants positifs
$\widetilde x,\widetilde y$, le franchissement du revêtement est
\[
 b(x,y)=\frac{\widetilde x+\nu_d-\widetilde y}{9}\in\mathbb Z^2.
\]
On a $b(y,x)=-b(x,y)$. Dans la catégorie libre des chemins
enregistrés, définissons le lecteur additif
\begin{equation}
 \mathfrak m(\gamma)=
 \left(\sum_{t=0}^{n-1}b(x_t,x_{t+1}),\ C_q(\gamma),\ n\right).
 \label{bim09:memoria-aditiva}
\end{equation}
La concaténation satisfait
$\mathfrak m(\gamma_2\circ\gamma_1)=
\mathfrak m(\gamma_1)+\mathfrak m(\gamma_2)$.
C'est un cocycle additif sur les trajectoires enregistrées : le registre conserve
les arêtes d'aller et de retour comme des événements différents.

Par \eqref{bim09:frontera-inversion},
\begin{equation}
 \mathfrak m(\gamma^\dagger\circ\gamma)=
 \bigl(0,\ 2C_q(\gamma)+q(x_n)-q(x_0),\ 2n\bigr),
 \label{bim09:memoria-retorno}
\end{equation}
où $\gamma^\dagger\circ\gamma$ désigne la concaténation temporelle
de l'aller suivi du retour. La somme orientée des franchissements s'annule ;
l'occupation arithmétique et la longueur demeurent. Ce lecteur ne
passe pas au quotient qui identifie un parcours suivi de son inverse
à la trajectoire vide, puisque ce quotient élimine les événements
conservés par les deuxième et troisième coordonnées.

Il existe également une lecture paire purement frontalière :
\[
 V_b(\gamma)=\sum_{t=0}^{n-1}
 b(x_t,x_{t+1})b(x_t,x_{t+1})^{\mathsf T},\qquad
 V_b(\gamma^\dagger\circ\gamma)=2V_b(\gamma).
\]
La lecture impaire et la lecture paire enregistrent des propriétés distinctes des
mêmes arêtes ; toutes deux sont calculées avant l'oubli du chemin.

\paragraph{Application exacte à la fibre de signature \texorpdfstring{$555555$}{555555}.}

L'énumération finie de $\mathscr X_\times$ contient $324$ germes,
$26$ signatures et $24$ germes de signature $555555$. La partition stable
sous $P,H$ raffine cette fibre en trois classes de huit éléments. Elles peuvent
être étiquetées par leur signature après une avancée :
\[
 \mathscr C_{177},\quad\mathscr C_{717},\quad\mathscr C_{771}.
\]
L'action de demi-tour est
\[
 H\mathscr C_{177}=\mathscr C_{177},\qquad
 H\mathscr C_{717}=\mathscr C_{771},\qquad
 H\mathscr C_{771}=\mathscr C_{717}.
\]
Le certificat fini joint vérifie ces identités sur tous les
représentants. Dans les vingt-quatre germes, les six accumulateurs
bruts de $f_\ell$ valent exactement cinq. Les quotients
décimaux $c_j^{10}$ s'annulent donc dans cette fibre. L'information supplémentaire
s'observe dans les positions, les évaluations entières et les quotients
arithmétiques, dont les lectures diffèrent de cette signature.

La coordonnée transversale fixe $a$ du parcours cardinal vaut $4$ ou
$5$ dans ces germes. Pendant un tour de neuf avancées, chaque
valeur de la coordonnée mobile $u\in D_9$ est parcourue une fois. Puisque
$\gcd(a,9)=1$, les résidus positifs $r(au)$ permutent $D_9$.
On en déduit, sans utiliser de constante physique,
\[
 \sum_{u=1}^9au=45a,\qquad
 \sum_{u=1}^9r(au)=45,\qquad
 \sum_{u=1}^9q(au)=5(a-1).
\]
Le tour suivant parcourt les mêmes positions en sens opposé.
Dans les dix-huit activations,
\begin{equation}
 \sum N=90a,\qquad\sum r=90,\qquad
 \sum q=10(a-1)=
 \begin{cases}30,&a=4,\\40,&a=5.\end{cases}
 \label{bim09:ocupacion-30-40}
\end{equation}
Le produit entier accumulé vaut respectivement $360$ ou $450$.
L'enroulement net est nul ; la somme des quotients et les dix-huit
incidences conservent l'histoire de la lecture. Chacune des trois
classes de signature contient des représentants des deux valeurs transversales.
La même signature et la même classe de transport de ce quotient fini
peuvent donc coexister avec des occupations arithmétiques distinctes.

Ces résultats relient la mémoire conservée à des opérations
concrètes : permutation des demi-cycles, différences de frontière,
quotients arithmétiques et registre d'incidences. L'identification
physico-mathématique de tout lecteur ultérieur conserve ses applications
propres ; les identités démontrées ici appartiennent à la dynamique
interne du canal produit.
